\section{量子K理论（Xiaohan Yan 与 You-Cheng Chou）}%
\label{sec:Quantum K-theory (Xiaohan Yan and You-Cheng Chou)}

\subsection{概述}%
\label{sub:Overview}

格罗莫夫--威滕不变量是某种在光滑射影簇 \(X\) 中“计数”曲线的方法。GW 理论给出形如
\[ \ab<\alpha_1, \ldots, \alpha_n>_{g, n, \beta} \in \Q \]
的数，其中 \(\alpha_i \in H^*(X)\) 是各类，\(g\) 是亏格，\(n\) 是标记点的个数，\(\beta \in H_2(X, \Z)\) 是曲线类。这已被用于研究镜像对称，并且在镜像对称之外还给出了双有理不变量。我们现在的目标是把这整套理论提升到 K理论。量子K理论给出数
\[ \ab<V_1, \ldots, V_n>_{g, n, \beta} \in \Z \]
其中 \(V_1, \ldots, V_n \in K(X)\)。它们满足各种差分方程\footnote{与上同调中的微分方程相比较。}，并在 3D 镜像对称与表示论中有诸多应用。

\subsection{K理论}%
\label{sub:K-theory}

\begin{defn}
    设 \(X\) 是光滑射影簇。则 \(X\) 的 \textit{K理论} \(K(X)\) 是 \(X\) 上向量丛的 Grothendieck 群，满足关系：若
    \[ 0 \to V_1 \to V_2 \to V_3 \to 0 \]
    是短正合列，则 \(V_2 = V_1 + V_3\)。
\end{defn}

\begin{rmk}\leavevmode
    \begin{enumerate}
        \item \(K(X)\) 是以直和与张量积为运算的环；
        \item 陈特征标给出同构 \(K(X) \to A^*(X)\)，至少在 \(\Q\) 上成立。
    \end{enumerate}
\end{rmk}

\begin{exm}
    设 \(X = \P^n\) 且 \(P = \mc{O}(-1)\)。则我们有
    \[ K(X) = \C[P^{\pm 1}]/(P-1)^{n+1} \xrightarrow{\on{ch}} A^*(X) = \C[H]/H^{n+1}. \]
\end{exm}

对于不假定光滑的 Deligne--Mumford 叠，有两个版本
\[ K_0(X) \supset K^0(X), \]
其中 \(K_0\) 表示凝聚层，\(K^0\) 表示完美复形。它们在光滑射影簇上一致，故从现在起 \(K(X)\) 均指 \(K_0(X)\)。存在张量积
\[ K^0(X) \otimes K_0(X) \to K_0(X) \]
以及沿真（proper）态射的推前
\[ f_* \colon K(X) \to K(Y) \qquad V \mapsto \sum (-1)^i R^i f_* V. \]
若 \(f\) 是完美态射（例如平坦、正则嵌入等），则存在拉回
\[ f^* \colon K(Y) \to K(X) \qquad W \mapsto \sum (-1)^i \on{Tor}_i^{f^{-1} \mc{O}_Y} (\mc{O}_X, f^{-1} W). \]

\subsection{曲线模空间}%
\label{sub:Moduli of curves}

回忆到 \(X\) 的稳定映射模空间
\[ \bar{M}_{g,n}(X, \beta) = \ab\{ f \colon (\Sigma, p_1, \ldots, p_n) \to X \mid g(\Sigma) = g, f_*[\Sigma] = \beta, \abs{\Aut f} < \infty \}. \]
它有求值映射
\[ \ev_i \colon \bar{M}_{g,n}(X, \beta) \to X \qquad (f, \Sigma, \ul{p}) \mapsto f(p_i). \]
还有万有余切线 \(L_i = \sigma_i^* \omega_{ \mc{C}/\bar{M} }\)。

该模空间是真 DM叠，具有完美障碍理论，从而给出虚结构层 \(\mc{O}^{\vir}_{g,n,\beta}\)。这来自\textit{虚切丛}的结构，其中
\[ \mathbb{T}_{(f, \Sigma, \ul{p})} \bar{M}_{g,n}(X, \beta) = \chi(\Sigma, f^* T_X) - \chi(\Sigma, T_{\Sigma}(-p_1 - \cdots - p_n)). \]
粗略地说，若 \(\bar{M}\) 是簇 \(Y\) 上向量丛 \(E\) 的截面的零点轨迹，则 \(\mc{O}^{\vir} = \mc{O}_Y \otimes \mr{Eu}(E)\)，其中 K 理论的欧拉类为 \(\mr{Eu}(E) = \sum (-1)^i \Lambda^i E^{\vee}\)。

\subsection{KGW 不变量}%
\label{sub:KGW invariants}

关于量子K理论，参见~\cite{Lee2004}；关于置换等变理论，参见~\cite{Givental2017K1,GiventalK7}。To\"en~\cite{Toen1999} 证明了下面用到的 Deligne--Mumford 叠上的 Riemann--Roch 公式。

\begin{defn}
    K 理论的 GW 不变量定义为
    \[ \ab<V_1, \ldots, V_n>_{g,n,\beta} \coloneqq \chi\ab(\bar{M}_{g,n}(X, \beta); \mc{O}^{\vir} \otimes \bigotimes_{i=1}^n \ev_i^* V_i). \]
\end{defn}

若 \(V_1, \ldots, V_n\) 是真正的向量丛，则该不变量是整数，但一般来说我们会考虑系数取自特征零域的 \(K_0\)。我们也可以定义引力后代
\[ \ab<V_1 L_1^{k_1}, \ldots, V_n L_n^{k_n}>_{g,n,\beta} \coloneqq \chi\ab(\bar{M}_{g,n}(X, \beta); \mc{O}^{\vir} \otimes \bigotimes_{i=1}^n ( \ev_i^* V_i )L_i^{k_i} ). \]
利用 \(S_n\) 通过置换标记点给出的在 \(\bar{M}_{g,n}(X, \beta)\) 上的作用，对象 \(\bigotimes_{i=1}^n \ev_i^* V\) 与 \(\bigotimes_{i=1}^n L_i^k\) 是 \(S_n\)-等变的。

\begin{defn}
    我们定义\textit{置换不变的 K 理论 GW 不变量}为
    \begin{align*}
        \ab<V L_1^k, \ldots, V L_n^k>_{g,n,\beta}^{S_n} &\coloneqq \chi\ab([ \bar{M}_{g,n}(X, \beta)/S_n ]; \mc{O}^{\vir} \otimes \bigotimes_{i=1}^n ( \ev_i^* V ) L_i^k) \\
        &= \dim \chi_{S_n} \ab(\bar{M}_{g,n}(X, \beta); \mc{O}^{\vir} \otimes \bigotimes_{i=1}^n ( \ev_i^* V ) L_i^k )^{S_n}.
    \end{align*}
\end{defn}

给定 \(t(q) = \sum_{n \in \Z} t_n q^n \in K(X)[q,q^{-1}]\)，我们可以定义
\[ \ab< t(L_1), \ldots, t(L_n)>^{S_n}_{g,n,\beta} \coloneqq \chi ([\bar{M}_{g,n}(X,\beta)/S_n], \cdots). \]
更一般地，若 \(H = S_{a_1} \times \cdots \times S_{a_m} \subset S_n\)，我们可以定义
\[ \ab<t_1(L_1), \ldots, t_1(L_{a_1}), t_2(L_{a_1+1}), \ldots, t_2(L_{a_1+a_2}), \ldots>_{g,n,\beta}^H. \]

\begin{rmk}
    我们想讨论与通常 GW 不变量的一些关系。回忆到 Hirzebruch--Riemann--Roch 指出
    \[ \chi(M; V) = \int_M \on{ch}(V) \on{td}(T_M). \]
    当 \(M\) 是光滑 DM叠时，T\"oen 证明了公式
    \[ \chi(M; V) = \int_{IM} \wt{\on{ch}}(V) \wt{\on{td}}(T_M), \]
    其中 \(IM = M \times_{\Delta, M \times M, \Delta} M\) 是 \(M\) 的惯性叠，波浪号表示取 orbifold 陈特征标与 Todd 类。 
\end{rmk}

\begin{exm}
    若 \(M = B\mu_r\)，则 \(IM = \bigsqcup_{g \in \mu_r} B \mu_r\)。若 \(M = \P(1, r)\) 是加权 \(\P^1\)，则 \(IM = M \sqcup \bigsqcup_{g \neq 1 \in \mu_r} B\mu_r\)。    
\end{exm}

\begin{exm}
    设 \(M = BG\)，其中 \(G\) 是有限阿贝尔群，且 \(V\) 是 \(BG\) 上的向量丛，这等同于 \(G\) 的一个表示。则 \(\chi(M, V) = \dim V^G\)。另一方面，我们有
    \[ \int_{IM} \wt{\on{ch}}(V) \cdot \wt{\on{td}}(T_M) = \int_{IM} \on{tr}_g V = \sum_{g \in G} \frac{1}{\abs{G}} \on{tr}_g V = \dim V^G, \]
    于是我们在此情形下验证了 Riemann-Roch。
\end{exm}

\subsection{公理化性质}%
\label{sub:Axiomatic properties}

\begin{lem}[{\cite{GiventalK7}}]
    我们有\textit{弦方程}
    \begin{align*}
        \ab< 1, x_1(L), \ldots, x_k(L), t, \ldots, t >_{g, 1+k+n, \beta}^{S_n} = \ab<x_1(L), \ldots, x_k(L), t^{\otimes n}>^{S_n}_{g,k+n,\beta} \\
        + \sum_i \ab< \cdots, \frac{x_i(L) - x_i(1)}{L-1}, \ldots, t^{\otimes n}>^{S_n}_{g,k+n,\beta}. 
    \end{align*}
\end{lem}

\begin{proof}
    证明与 GW 理论中相同，来自沿映射 \(\pi \colon\bar{M}_{g,n+1}(X,\beta) \to \bar{M}_{g,n}(X,\beta)\) 计算 \(\pi^* L_i\)。
\end{proof}

\begin{lem}[{\cite{GiventalK7}}]
    我们有\textit{膨胀子方程}
    \begin{align*}
        \ab<L-1, x_1(L), \ldots, x_k(L), t^{\otimes n}>_{g,1+k+n,\beta}^{S_n} ={}& (k-2)\ab<x_1(L), \ldots, x_k(L), t^{\otimes n}>_{g,k+n,\beta}^{S_n} \\
        &+ \ab<x_1(L), \ldots, x_k(L), t, t^{\otimes n-1}>_{g,k+n,\beta}^{S_{n-1}}.
    \end{align*}
\end{lem}

\begin{proof}
    我们利用
    \[ \pi_* L = k + \on{Ind}_{S_{n-1}}^{S_n}(1) - 1, \]
    其中 \(\pi\) 遗忘第一个标记点，\(L = \omega_{\pi} (p_1 + \cdots + p_n)\)。这来自存在留数映射
    \[ H^0(L) \xrightarrow{\on{Res}} \bigoplus_{i=1}^{k+n} \C p_i \xrightarrow{\sum} \C. \]
    中间项是 \(k\) 个平凡表示与一个标准表示的直和。
\end{proof}

\begin{notn}
    固定 \(\tau, t \in K(X)\)。我们将记
    \[ \ab<\!\ab<x_1(L_1), \ldots, x_k(L_k)>\!>_{g,k} \coloneqq \sum_{m,n,\beta} \frac{Q^{\beta}}{m!} \ab<x_1(L_1), \ldots, x_k(L_k), \tau^{\otimes m}, t^{\otimes n}>_{g,k+m+n,\beta}^{S_n}. \]
    我们令 \(\{\phi_{\alpha}\}\) 为 \(K(X)\) 的一组基，\(\{\phi^{\alpha}\}\) 为对偶基。然后定义配对
    \[ g_{\alpha\beta} = \chi(X, \phi_{\alpha} \otimes \phi_{\beta}). \]
    这并不完全是正确的配对，故我们考虑
    \[ G_{\alpha\beta} = g_{\alpha\beta} + \ab<\!\ab<\alpha,\beta>\!>_{0,2}. \]
    我们令 \(g^{\alpha\beta} = (g_{\alpha\beta})^{-1}\)，\(G^{\alpha\beta} = (G_{\alpha\beta})^{-1}\)。
\end{notn}

\begin{lem}[{\cite{Givental2000K}}]
    给定 \(A,B,C,D \in K(X)[q,q^{-1}]\)，我们有\textit{WDVV 方程}
    \begin{align*}
        \sum_{\alpha,\beta} \ab<\!\ab< A(L), B(L), \phi_{\alpha}>\! >_{0,3}& G^{\alpha\beta}\ab<\!\ab< C(L), D(L), \phi_{\beta}>\! >_{0,3} \\
        &= \sum_{\alpha,\beta} \ab<\!\ab< A(L), C(L), \phi_{\alpha}>\! >_{0,3} G^{\alpha\beta}\ab<\!\ab< B(L), D(L), \phi_{\beta}>\! >_{0,3}. 
    \end{align*}
\end{lem}

\begin{proof}
    拉回 \(\bar{M}_{0,4} = \P^1\) 上边界除子的等式。在 K理论中，当 \(D = \bigcup C_i\) 时，把 \(\mc{O}_D\) 与 \(\mc{O}_{C_i}\) 联系起来需要应用容斥原理，这解释了 \(G^{\alpha\beta}\) 的出现。
\end{proof}

\subsection{环路空间形式}%
\label{sub:Loop space formalism}

\begin{defn}
    \textit{环路空间} \((\mc{K}, \Omega)\) 由
    \[ \mc{K} = K(X)(q) \ps{Q}\]
    组成，带有辛形式
    \[ \Omega(f, g) = - \on{Res}_{q = 0, \infty} \ab<f(q^{-1}), g(q)>_X \frac{\d{q}}{q}. \]
\end{defn}

\begin{rmk}
    存在极化 \(\mc{K} = \mc{K}_+ \oplus \mc{K}_-\)，其中
    \[ \mc{K}_+ = K(X)[q, q^{-1}]\ps{Q} \]
    以及
    \[ \mc{K}_- = \ab\{ \frac{f}{g} \mid f, g \in K(X)\ps{Q}[q], g(0) \neq 0, \deg f < \deg g\}. \]
    例如，\(\frac{1}{1-q} \in \mc{K}_-\)。
\end{rmk}

\begin{defn}
    \textit{大 J-函数} \(J \colon \mc{K}_+ \to \mc{K}\) 由
    \[ \mathbf{t} = \mathbf{t}(q) \mapsto 1 - q + \mathbf{t}(q) + \sum_{\beta, \alpha} Q^{\beta} \phi^{\alpha} \ab< \frac{\phi_{\alpha}}{1-qL}, \mathbf{t}(L)^{\otimes n}>_{0, 1+n, \beta}^{S_n}. \]
    给出。
\end{defn}

\begin{rmk}
    注意到应用 Riemann-Roch，我们有
    \begin{align*}
        \chi\ab([\bar{M}/S_n], \frac{1}{1-qL}) &= \int_{I[\bar{M}/S_n]} \on{ch}\ab(\frac{1}{1-qL}) \on{td}.
    \end{align*}
    每一项都变成 \(\on{ch}\ab(\on{tr}_g \frac{1}{1-qL})\) 的形式，这给出形如 \(\on{ch}\ab(\frac{1}{1-q\zeta L})\) 的东西，其中 \(\zeta\) 是某个单位根。于是我们可以展开
    \[ \frac{1}{1-q\zeta L} = \frac{1}{(1-q \zeta) - q\zeta (L-1)} = \sum_k \frac{(q\zeta)^k}{(1-q\zeta)^{k+1}}(L-1)^k. \]
\end{rmk}

\begin{defn}
    \textit{Givental 锥}定义为 \(\mc{L} \coloneqq \Ima J \subset \mc{K}\)，其中我们考虑在 \(Q\)-进拓扑下 \(1-q\) 附近的芽。
\end{defn}

\begin{exm}
    设 \(X = \P^n\)。则我们看到
    \[ J(0) = (1-q) \sum_{d \geq 0} \frac{Q^d}{\prod_{\ell=1}^d (1-Pq^{\ell})^{n+1}}, \]
    其中 \(P = \mc{O}(-1)\)。如前，我们可以展开
    \[ \frac{1}{1-P q^{\ell}} = \frac{1}{(1-q^{\ell}) - q^{\ell}(P-1)} = \sum_{k=0}^{\infty} \frac{q^{\ell k}}{(1-q^{\ell})^{k+1}}(P-1)^k. \]
\end{exm}

\begin{prop}[{\cite{GiventalK7}}]
    事实上，\(\mc{L}\) 确实是 \(\mc{K}\) 中的一个锥，并且我们有
    \[ \mc{L} = \bigcup_{t \in K(X)\ps{Q}} (1-q) S_{t}^{-1}(q) \mc{K}_+. \]
\end{prop}

我们还可以引入一个“巨大”的 J-函数
\[ J(\mathbf{t}, \tau) = 1-q + \mathbf{t}(q) + \tau + \sum_{\beta, k, \alpha} \frac{Q^{\beta}}{k!} \phi^{\alpha} \ab< \frac{\phi_{\alpha}}{1-qL}, \mathbf{t}(L)^{\otimes n}, \tau^{\otimes k}>_{0,1+n+k,\beta}^{S_n}. \]
若固定 \(\mathbf{t}(q)\)，则 \(J(\tau)\) 来自 \(\tau\) 方向 Frobenius 流形上一个\textit{微分方程}的基本解。此外，\(\mc{L}^{\mathbf{t}}=  \Ima J(\tau)\) 实际上是一个拉格朗日锥。

\begin{exm}
    考虑由 Adams 运算给出的 \(f \colon \C[\Lambda_1^{\pm}, \Lambda_2^{\pm}] \to \C[\Lambda_1^{\pm}, \Lambda_2^{\pm}]\)
    \[ t \mapsto \on{tr}_{(1\cdots n)} t^{\otimes n} = \Psi^n (t). \]
    因此，它把 \(\Lambda_1 \mapsto \Lambda_1^n\)、\(\Lambda_2 \mapsto \Lambda_2^n\)。然而，如何在 \(t\) 中求导并不清楚，因为我们需要 \(nv t^{n-1} = \Psi^n(v)\)，所以 J-函数不满足关于 \(\mathbf{t}\) 的量子微分方程。
\end{exm}

\subsection{等变 K理论与局部化}%
\label{sub:Equivariant K-theory and localization}

关于不动点递推与射影空间的计算，参见~\cite{GiventalK2}。

设环面 \(T = (\C^{\times})^r\) 作用在光滑射影簇 \(X\) 上。

\begin{defn}
    \textit{等变 K理论} \(K_T(X)\) 是 \(X\) 上 \(T\)-等变向量丛的 Grothendieck 群。
\end{defn}

\begin{exm}
    例如，我们有
    \[ K_T(\on{pt}) = \on{Rep}(T) = \C[\Lambda_1^{\pm}, \ldots, \Lambda_r^{\pm}], \]
    其中 \(\Lambda_i\) 是权为 \((0, \ldots, 1, \ldots, 0)\) 的一维表示。
\end{exm}

\begin{rmk}
    对任意 \(X\)，\(K_T(X)\) 是 \(K_T(\mr{pt})\)-代数。
\end{rmk}

\begin{exm}
    考虑 \(\P^n\) 的完全等变 K理论。则我们有
    \[ K_T(\P^n) = \C[P^{\pm}, \Lambda_0^{\pm}, \ldots, \Lambda_n^{\pm}] \big/ \prod_{i=0}^n (P-\Lambda_i). \]
    当我们特殊化 \(\Lambda_i = 1\) 时，就恢复出非等变 K理论。
\end{exm}

\begin{defn}
    我们可以把等变欧拉示性数定义为推前
    \[ \chi_T = \pi_* \colon K_T(X) \to K_T(\mr{pt}). \]
\end{defn}

\begin{exm}
    例如，我们有 \(\chi_T(\mc{O}_{\P^1}(1)) = \Lambda_0^{-1} + \Lambda_1^{-1}\)。类似地，我们有
    \[ \chi_T(\mc{O}(2)) = \Lambda_0^{-2} + \Lambda_0^{-1} \Lambda_1^{-1} + \Lambda_1^{-2}. \]
\end{exm}

\begin{notn}
    注意到 \(\on{Frac}(T) = \on{Frac}(\on{Rep}(T))\)。于是我们令
    \[ K_T^{\loc}(X) \coloneqq \on{Frac}(T) \otimes_{\on{Rep}(T)} K_T(X). \]
\end{notn}

我们需要以下事实。设 \(F\) 是 \(T\) 在 \(X\) 上的不动点集，并考虑 \(\iota \colon F \hookrightarrow X\)。则存在 \(\on{Frac}(T)\)-代数的同构
\[ K_T^{\loc}(X) \xrightarrow{\iota^*} K_T^{\loc}(F). \]

\begin{rmk}
    逆由
    \[ (\iota^*)^{-1} = \bigoplus_{i=1}^r \frac{\iota_{i*}}{\on{Eu}(N_{F_i/X})}, \]
    给出，其中 \(\on{Eu}(V) \coloneqq \sum_{j = 0}^{\on{rk} V} (-1)^j \Lambda^j V^*\) 是 K 理论的欧拉类。
\end{rmk}

\begin{exm}
    在 \(\P^n\) 上，我们有 \(n+1\) 个孤立不动点 \(F_i\)，于是得到
    \[ \frac{\iota_{i*} \mc{O}_{F_i}}{\on{Eu}(N_{F_i/\P^n})} = \prod_{j \neq i} \frac{P - \Lambda_j}{\Lambda_i - \Lambda_j}. \]
    它们是 \(K_T^{\loc}(\P^n)\) 的幂等元。
\end{exm}

利用等变局部化，我们得到
\[ \chi_T(X; V) = \sum_i \chi\ab(F_i; \frac{\iota_i^* V}{\on{Eu}(N_{F_i/X})}). \]

\begin{rmk}
    在 \(\Lambda = 1\) 的极限下，我们可以恢复非等变欧拉示性数。
\end{rmk}

\begin{exm}
    设 \(X = \P^1\)。则我们计算得
    \begin{align*}
        \chi_T(\P^1, \mc{O}(1)) &= \chi\ab(F_0, \frac{\mc{O}(1)|_{F_0}}{\on{Eu}(T_{F_0/\P^1})}) + \chi\ab(F_1, \frac{\mc{O}(1)|_{F_1}}{\on{Eu}(T_{F_1/\P^1})}) \\
        &= \frac{\Lambda_0^{-1}}{1 - \frac{\Lambda_0}{\Lambda_1}} + \frac{\Lambda_1^{-1}}{1 - \frac{\Lambda_1}{\Lambda_0}} \\
        &= \Lambda_0^{-1} + \Lambda_1^{-1}. 
    \end{align*}
\end{exm}

从现在起，我们假设 \(X\) 是 GKM 流形。我们现在的目标是在 \(\bar{M}_{g,n}(X, \beta)\) 上使用局部化来研究 \(X\) 的大 J-函数。任意 \(T\)-不动稳定映射 \(f \colon (\Sigma, \ul{p}) \to X\) 有两种不可约分支：完全映到不动点的分支，以及通过一维轨道的分歧覆盖来映的分支。遗憾的是，我们不能通过虚局部化计算大 J-函数，但仍可以研究它。

设 \(f = f(q) \in \mc{K} = K_T^{\loc}(X)(q)\ps{Q}\)。把不动点记为 \(F_{\alpha}\)，\(\alpha = 1, \ldots, r\)，含入为 \(\iota_{\alpha} \colon F_{\alpha} \hookrightarrow X\)。然后记
\[ f_{\alpha}(q) = \iota_{\alpha}^* f(q) \in K_T^{\loc}(\on{pt})(q)\ps{Q}. \]
则 \(f \in \mc{L}\) 当且仅当对所有 \(\alpha\)，以下条件成立：
\begin{enumerate}
    \item 给定 \(\lambda \in T_{F_{\alpha}} X\) 和 \(m \in \N_+\)，\(f_{\alpha}(q)\) 在 \(q = \lambda^{1/m}\) 处至多有单极点，并且
        \[ \on{Res}_{q = \lambda^{1/m}} f_{\alpha}(q) \frac{\d{q}}{q} = -\frac{Q^m}{m} \cdot \frac{\on{Eu}(T_{F_{\alpha}/X})}{\on{Eu}(T^{\vir}_{\phi_{\alpha,\alpha',m}}\bar{M}_{0,2}(X, m \cdot \beta(\alpha, \alpha')))} \cdot f_{\alpha'}(\lambda^{1/m}). \]
        这里，\(\lambda \in T_{F_{\alpha}}X\) 意味着 \(\lambda\) 是 \(T_{F_{\alpha}} X\) 的一维子表示。此外，\(\beta(\alpha,\alpha')\) 是由权 \(\lambda\) 确定的 \(T\) 的一维轨道的次数，\(\phi_{\alpha, \alpha', m}\) 参数化该一维轨道的 \(m:1\) 覆盖。
    \item \(f_{\alpha}(q)\) 的其余极点位于 \(0, \infty\) 和单位根处，在其他地方正则。此外，还有各种次要条件。
\end{enumerate}

\begin{exm}
    设 \(X = \P^n\)。则
    \[ J_T^X(0) = (1-q) \sum_{d \geq 0} \frac{Q^d}{\prod_{i=0}^n \prod_{\ell = 1}^d \ab(1 - \frac{P}{\Lambda_i} q^{\ell})} \in K_T(X)(q) \ps{Q}. \]
    于是这推出
    \[ J_T^X(0) |_{F_j} = (1-q) \sum_{d \geq 0} \frac{Q^d}{\prod_{i=0}^n \prod_{\ell = 1}^d \ab(1 - \frac{\Lambda_j}{\Lambda_i}q^{\ell})}, \]
    故极点在 \(q = \ab(\frac{\Lambda_i}{\Lambda_j})^{1/\ell}\) 处。还要注意到
    \[ T_{F_j} X = \frac{\Lambda_0}{\Lambda_j} + \cdots + \frac{\Lambda_n}{\Lambda_j} - 1. \]
\end{exm}

\subsection{量子K/GV 对应}%
\label{sub:Quantum K/GV correspondence}

五次对应在~\cite{GaroufalidisScheidegger2022} 中猜想，并在~\cite{ChouLeeQuintic} 中证明。下面的 Calabi--Yau 三维簇公式归功于 Chou--Lee~\cite{ChouLee2025}。

我们这里的目标是在 Calabi-Yau 三维簇上计算 \(J^K(0)\)。以下结果由 Jockers-Mayr 和 Garoufalidis-Scheidegger 猜想。

\begin{thm}[{\cite{ChouLee2025}}]
    我们有
    \begin{align*}
        J^K(0) &\coloneqq (1-q) + \sum_{\beta, \alpha} Q^{\beta} \Phi^{\alpha} \ab<\frac{\Phi_{\alpha}}{1-qL}>_{0,1,\beta}^{X,K} \\
        &= (1-q) + \sum_{\beta} \sum_{r=1}^{\infty} \mr{GV}_{\beta} Q^{r\beta} \ab[\sum_{j=1}^{n_1} \Phi^{1j} a(r, q^r) + \Phi^{01} b(r, q^r)],
    \end{align*}
    其中 \(\ab\{\Phi_{\alpha}\}_{\alpha=1}^N\) 是 \(K^0(X)\) 的一组基，分解为
    \[ \ab\{\Phi_{\alpha}\}_{\alpha=1}^N = \bigsqcup_{i=0}^3 \ab\{\Phi_{ij}\}_{j=1}^{n_i} \]
    使得对固定的 \(i\)，集合 \(\ch(\Phi_{ij})\) 构成 \(H^{2i}(X)\) 的一组基。然后我们选取对偶基 \(\ab\{\Phi^{\alpha}\}_{\alpha} = \ab\{\Phi^{ij}\}_{i,j}\)。此外，我们令多重覆盖公式
    \begin{align*}
        a(r, q^r) &= \frac{r-1}{1-q^r} + \frac{1}{(1-q^r)^2}, \\
        b(r, q^r) &= \frac{r^2-1}{1-q^r} + \frac{3}{(1-q^r)^2} - \frac{2}{(1-q^r)^3}.
    \end{align*}
    最后，GV 表示 Gopakumar-Vafa 不变量，它在亏格零时由
    \[ \mr{GW}_{0,\beta} = \sum_{r \mid \beta} \frac{1}{r^3} \mr{GV}_{\beta/r}. \]
    定义。由 M\"obius 反演，我们看到
    \[ \mr{GV}_{\beta} = \sum_{r \mid \beta} \frac{\mu(r)}{r^3} \mr{GW}_{0,\beta/r}. \]
\end{thm}

\begin{cor}[{\cite{ChouLee2025}}]
    我们有
    \[ \ab<1>_{0,1,\beta}^K = \sum_{r \mid \beta} r^2 \mr{GV}_{\beta/r}. \]
    应用 M\"obius 反演，我们得到\footnote{注意到这一结果由 Ionel--Parker~\cite{IonelParker2018} 在所有亏格下得到。}
    \[ \mr{GV}_{\beta} = \sum_{r \mid \beta} r^2 \mu(r) \ab<1>_{0,1,\beta/r}^K \in \Z. \]
    此外，我们可以通过\footnote{这就好比若我们应用 Hirzebruch-Riemann-Roch 来计算 K 理论配对，最高次项只看到两个输入项的陈特征标。}
    \[ \lim_{q \to 1} (1-q)^3 \ab<\frac{1}{1-qL}>_{0,1,\beta}^{X,K} = \sum_{r \mid \beta} \frac{-2}{r^3} \mr{GV}_{\beta/r} = -2 \mr{GW}_{\beta} = \ab<\psi>_{0,1,\beta}^{\mr{GW}}. \]
    恢复通常的 GW 不变量。
\end{cor}

粗略地，我们可以把 \(J^K(0)\) 视为 \(J^K(0) \in \Z\ps{q}\ps{Q}\)，但更精确地说，我们可以考虑
\[ J^K(0) \in 1-q + \mc{K}_-. \]
这里，我们展开 \(\frac{1}{1-qL} = \sum q^i L^i\)。 
我们还需要考虑这样一个事实：若 \(M\) 是 DM叠，\(L\) 是 \(M\) 上的线丛，则函数
\[ f^{\Phi}(q) \coloneqq \chi\ab(M, \frac{\Phi}{1-qL}) \in \Z \ps{q} \]
在 \(\abs{q} < \ep\) 时收敛到有理函数。注意到极点只在单位根处。

这一事实可以用 orbifold Riemann-Roch 定理证明。把惯性叠分解为
\begin{align*}
    IM &= \bigsqcup_{i \in I} M_i \\
    &= \bigsqcup_{\zeta} M_{\zeta},
\end{align*}
其中 \(M_{\zeta}\) 是使得 \(g\) 以特征值 \(\zeta\) 作用在 \(L\) 上的扭曲扇区之并。利用这一点，我们可以计算
\begin{align*}
    \chi\ab(M, \frac{1}{1-qL}) &= \sum_{i \in I} \int_{M_i} \on{td}(M_i) \on{ch} \ab(\frac{ \on{Tr}_g \ab(\frac{1}{1-qL}) }{\lambda_{-1}(N_{M_i/M}^*)}) \\
    &= \sum_{\zeta} \int_{M_{\zeta}} \on{td}(M_{\zeta}) \sum_{k \geq 0} \frac{(\zeta q)^k}{(1-\zeta q)^{k+1}} \frac{( \on{ch}(L)-1 )^k}{\ch \lambda_{-1}(N^*)},
\end{align*}
因此所有极点都必须在单位根处。这里，我们使用展开
\[ \frac{1}{1-\zeta q L} = \sum_{k \geq 0} \frac{(\zeta q)^k}{(1-\zeta q)^{k+1}} (L-1)^k \]
并取陈特征标。

\begin{exm}
    设 \(L = \mc{O}(1)\) 在 \(\P^1\) 上。则我们有
    \begin{align*}
        \chi\ab(\P^1, \frac{1}{1-qL}) &= \chi\ab(\P^1, \frac{1}{1-q} + \frac{q}{(1-q)^2}(L-1)) \\
        &= \sum_{i \geq 0} q^i (i+1) \\
        &= \frac{1}{(1-q)^2}.
    \end{align*}
\end{exm}

\begin{exm}
    我们现在考虑加权射影空间。则我们有
    \begin{align*}
        \chi\ab(\P(w_0, \ldots, w_n), \frac{1}{1-qL}) &= \prod_{i=0}^n \frac{1}{1-q^{w_i}}. 
    \end{align*}
\end{exm}

